% Test-case specification TC-0014 (CTest test TC-0014-x87math).
\documentclass[letterpaper]{article}

\usepackage[letterpaper,margin=1in]{geometry}
\emergencystretch=3em

% Compile from the repository root so this path resolves to
% the pinned WSP submodule.
\input{wsp/testing/test-case-library.tex}

\begin{document}

\def\TCID{TC-0014}
\def\TCTitle{%
  The x87 math helpers and the wide multiply}
\def\TCPurpose{%
  Verify \texttt{\_\allowbreak{}CIsin} and the other
  \texttt{\_\allowbreak{}CI*} helpers, the
  \texttt{\_\allowbreak{}adj\_\allowbreak{}*} divides and
  \texttt{\_\allowbreak{}allmul}.}
\def\TCPriority{\TCPriorityHigh}

\def\TCRequirementRef{%
  REQ-0006}
\def\TCDesignRef{%
  Not applicable.}
\def\TCFeatureRef{%
  \texttt{tests/\allowbreak{}vbrun\_\allowbreak{}x87math\_\allowbreak{}tests.\allowbreak{}cpp},
  \texttt{src/\allowbreak{}helpers\_\allowbreak{}x87math.\allowbreak{}cpp}}
\def\TCTestTechnique{%
  Back-to-back comparison; boundary values.}

\def\TCPreconditions{%
  The library and the test program are built
  (\texttt{scripts/\allowbreak{}Invoke-Build.\allowbreak{}ps1}) and the test
  is registered with CTest.}
\def\TCEnvironment{%
  Windows with the toolchain of
  \texttt{scripts/\allowbreak{}Invoke-Build.\allowbreak{}ps1}; the local
  gate builds and runs x64 Debug, CI builds and runs the other
  architectures. The library is built from this repository in
  \texttt{out/\allowbreak{}build}.}
\def\TCAssumptions{%
  Not applicable.}

\def\TCInputData{%
  Operands of each helper, including the cases that round or raise.}
\def\TCInitialState{%
  No state from earlier runs is needed; the test makes what it needs.}

\def\TCProcedure{%
  \begin{enumerate}
    \item CTest runs \texttt{vbrun\_\allowbreak{}x87math\_\allowbreak{}tests}.
    \item On 32-bit x86, call the reference's sequences with the operand on the
      x87 stack; elsewhere call the exports that take doubles.
    \item Compare each result with the expected value within the tolerance of
      the check.
    \item CTest records the exit code and, on failure, the program's output
      (\texttt{--output-on-failure}).
  \end{enumerate}}
\def\TCExpectedResult{%
  All checks pass and the program exits 0.}
\def\TCPostConditions{%
  No file or registry state remains that later tests depend on.}

\TCTemplate

\end{document}
